\chapter{תפעול ותוכנית ביצוע}\label{ch:operations}

התפעול הוא היתרון התחרותי המעשי של גגון. גג אחד הוא מתקן קטן שחוזר על עצמו, ולכן כל שיפור בתהליך מוכפל ב-\RRroofsFinal{} גגות. פרק זה מתאר כיצד גג נכנס לפעילות, איך הוא מנוהל מדי יום ואיך התוצרת מגיעה ללקוח.

\section{הצטרפות גג: משלב הסקר עד הקטיף הראשון}\label{sec:onboarding}

תהליך ההתקנה נמשך כ-6 שבועות מרגע חתימת החוזה, ומורכב מארבעה שלבים.

\begin{enumerate}
  \item \textbf{סקר מבני}: מהנדס מאשר שהגג עומד בעומס של המערכת, מים וצמחים, ובודק גישה, ניקוז וחיבור למים.
  \item \textbf{היתרים ואישורים}: בקשה לרשות המקומית, אישור כיבוי אש ובטיחות, והסכמת הדיירים.
  \item \textbf{התקנה}: מערכת \en{NFT} עם חיישנים, משאבות ומערכת הצללה. ההתקנה מודולרית ונעשית ברובה בימים בודדים.
  \item \textbf{הרצה}: שתילה, כיול חיישנים ב-\en{RootOS} וקטיף ראשון.
\end{enumerate}

תקופת ההתקנה והאישורים היא הנקודה הרגישה ביותר בציר הזמן, ולכן הצוות מקדים את הסקר המבני לשלב המכירה.

\section{תפעול יומי: \en{RootOS} ומגדלים אנושיים}\label{sec:daily}

\en{RootOS} קורא נתוני \en{pH}, מוליכות חשמלית (\en{EC}), טמפרטורה וספיקה, ומפעיל משאבות, הצללה והשקיה. מודל ראייה מזהה מחלות בעלים. פעולות חריגות מועברות למגדל אנושי, כדי שהמערכת לא תפעל לבדה במצבים לא צפויים. מגדל אחד מפקח על מספר גגות מרחוק ומגיע לגג לביקור שבועי.

כך נשמרת עלות עבודה נמוכה ללא פגיעה באיכות. עקרונות הגידול ההידרופוני מבוססים על ספרות מקצועית כגון \cite{resh2012}, והוא מותאם לתנאי הגג בישראל.

\section{קטיף והובלה אחרונה}\label{sec:logistics}

הקטיף נעשה בבוקר, והאספקה למסעדות ולתושבים בתוך שעות. הגג קרוב ללקוח, ולכן המסלול קצר. מסלול האספקה נבנה לפי אשכול גגות בשכונה אחת, כדי להקטין את הקילומטראז'. מסעדות מקבלות אספקה לפי לוח קבוע, ותושבים מקבלים סל בלובי הבניין או בדלת. כמות התוצרת השנתית בגג היא \RRroofKg{} ק"ג, ולכן ניתן לשרת את שני הערוצים מאותו גג.

\section{איכות ובטיחות מזון}\label{sec:quality}

גידול ללא קרקע מקטין סיכוני זיהום, אך מוסיף סיכון מים. אנו מנטרים איכות מים, מבצעים בדיקות מעבדה תקופתיות ושומרים רישום קטיף לפי גג ותאריך, כך שניתן לעקוב אחר כל אצווה. הצוות פועל לפי עקרונות \en{HACCP} ושואף לעמוד בדרישות משרד הבריאות והחקלאות לגידול ולשיווק ירקות טריים.

\begin{riskbox}[סיכון תפעולי]
זיהום מים או כשל מערכת השקיה עלולים להשפיע על יבול שלם ועל מוניטין הבניין. התגובה: חיישנים רציפים, גידול מקביל בכמה גגות והתראה אוטומטית למגדל.
\end{riskbox}

\section{לוח זמנים לביצוע}\label{sec:timeline}

איור~\ref{fig:timeline} מציג את אבני הדרך מ-2027 עד 2031: מ-\RRroofsFirst{} גגות ראשונים, דרך גיוס ה-\en{Series A}, ועד \RRroofsFinal{} גגות ורווחיות תפעולית (\en{EBITDA}) מ-\RRbreakEvenYear. הצמיחה מוגבלת על ידי קצב ההתקנה של צוות התפעול ולא על ידי ביקוש.

\begin{figure}[htbp]
  \centering
  \resizebox{\linewidth}{!}{% Five-year roadmap — milestones on a horizontal time axis (right = start, RTL).
\begin{tikzpicture}[x=-2.9cm,
  ms/.style={circle,fill=#1,draw=white,line width=1.2pt,minimum size=4mm,inner sep=0},
  lab/.style={font=\scriptsize,align=center,text width=2.6cm}]
  \draw[line width=4pt,rrForest!25,line cap=round] (0,0) -- (5,0);
  \foreach \i/\y in {0/2027,1/2028,2/2029,3/2030,4/2031,5/2032} {
    \draw[rrForest,line width=0.8pt] (\i,-0.12) -- (\i,0.12);
    \node[font=\sffamily\bfseries\small,text=rrForest] at (\i,-0.45) {\y};}
  \foreach \t/\c/\up/\an/\txt in {
      0.15/rrLeaf/1/south/{סבב \en{Seed}\\ \RRseed{} מיליון \shekel},
      0.6/rrForest/-1/north/{\RRroofsFirst{} גגות פיילוט\\ בגוש דן},
      0.85/rrLeaf/1/south/{\en{RootOS} גרסת\\ ייצור ראשונה},
      2.1/rrTerra/-1/north/{סבב \en{A}\\ \RRseriesA{} מיליון \shekel},
      2.6/rrLeaf/1/south/{לקוחות \en{SaaS}\\ ראשונים},
      3.3/rrForest/-1/north/{\en{EBITDA}\\ חיובי},
      4.4/rrTerra/1/south/{\RRroofsFinal{} גגות\\ ברחבי גוש דן}} {
    \node[ms=\c] at (\t,0) {};
    \draw[\c,line width=0.7pt] (\t,0) -- (\t,0.85*\up);
    \node[lab,anchor=\an] at (\t,0.85*\up) {\txt};}
\end{tikzpicture}
}
  \caption{מפת דרכים לשנים 2027--2031.}
  \label{fig:timeline}
\end{figure}

\begin{insightbox}[תובנה]
תהליך ההצטרפות של 6 שבועות הוא תהליך ניתן לחזרה. אחרי כמה גגות הוא הופך לתבנית, וקצב ההתקנה יכול לגדול בלי גידול דומה בצוות.
\end{insightbox}
